\begin{abwframe}
\node[anchor=north west,inner sep=0pt,outer sep=0pt] at (38.551,-150.803) {\begin{minipage}[t][363.546bp][t]{877.846bp}\setlength{\parindent}{0pt}\setlength{\parskip}{0pt}{\TimesNR\fontsize{32}{35.84}\selectfont\color{c000000}\raggedright \hspace*{10bp}\makebox[12bp][l]{\textbullet}Generative models aim to model the joint probability distribution of the input data and the target labels. \par}{\TimesNR\fontsize{32}{35.84}\selectfont\color{c000000}\raggedright \hspace*{10bp}\makebox[12bp][l]{\textbullet}In other words, they learn how data is generated from the underlying probability distribution.\par}{\TimesNR\fontsize{32}{35.84}\selectfont\color{c000000}\raggedright \hspace*{10bp}\makebox[12bp][l]{\textbullet}So, if a model learns about cat pictures, it can be asked to produce an original picture of a cat! \par}\end{minipage}};
\node[anchor=north west,inner sep=0pt,outer sep=0pt] at (38.551,-65.764) {\begin{minipage}[t][73.709bp][c]{878.677bp}\setlength{\parindent}{0pt}\setlength{\parskip}{0pt}{\TimesNR\fontsize{44}{49.28}\selectfont\color{cBC0031}\bfseries\raggedright Generative models\par}\end{minipage}};
\end{abwframe}
